% from ega4-20-fr.tex
\oldpage[IV-4]{244}
Lorsque $X$ vérifie l'une des conditions~a),~b) de \sref{IV.20.3.6}, alors on a vu \sref{IV.20.2.11} que les pseudo-fonctions sur $X$ s'identifient aux fonctions méromorphes sur $X$.
En vertu de \sref{IV.20.1.12} et de \sref{IV.20.2.13}[(i)], si l'on suppose seulement que le morphisme $f:X'\to X$ est \emph{plat}, alors, pour toute pseudo-fonction $\vphi$ sur $X$, $\vphi\circ f$ est définie et l'on a $\operatorname{dom}(\vphi\circ f)=f^{-1}(\operatorname{dom}(\vphi))$.
\end{remark}

\subsection{Propriétés des domaines de définition des fonctions rationnelles}
\EGAReaderTarget{ega4.c04.subsec.20-4}{IV.20.4}

\begin{env}[20.4.1]
\EGAReaderTarget{ega4.c04.statement.20-4-1}{IV.20.4.1}
Soient $X$, $Y$ deux $S$-préschémas, $\omega$ un pseudo-$S$-morphisme de $X$ dans $Y$.
Soit $u$ un $S$-morphisme $U\to Y$ appartenant à $\omega$, où $U$ est schématiquement dense dans $X$, et considérons le \emph{graphe} $\Gamma_u$, sous-préschéma de $U\times_S Y$ \sref[I]{I.5.3.11}, donc sous-préschéma de $X\times_S Y$.
Supposons que ce sous-préschéma admette une \emph{adhérence} \sref[I]{I.9.5.11} dans $X\times_S Y$ ; si $j:\Gamma_u\to X\times_S Y$ est l'injection canonique, ce sera le cas lorsque le $\sh{O}_{X\times_S Y}$-Module $j_*(\sh{O}_{\Gamma_u})$ est quasi-cohérent ; il résulte de la définition de la classe d'équivalence $\omega$ \sref{IV.20.2.1} que $j_*(\sh{O}_{\Gamma_u})$ ne dépend pas du représentant $u$ considéré, donc le préschéma adhérence de $\Gamma_u$ est alors bien déterminé par $\omega$ ; on dit que ce sous-préschéma fermé de $X\times_S Y$ est le \emph{graphe} du pseudo-$S$-morphisme $\omega$ et on le note $\Gamma_\omega$.
On notera que $\Gamma_\omega$ est défini s'il existe un morphisme $u:U\to Y$ de la classe $\omega$ tel que $U$ soit \emph{rétrocompact} dans $X$ et si de plus $Y$ est quasi-séparé sur $S$, car alors l'injection $j$ est un morphisme quasi-compact et quasi-séparé (\sref[I]{I.2.2} et \sref[I]{I.7.4}) ; la première hypothèse sera toujours vérifiée lorsque $X$ est \emph{localement noethérien}.
Notons aussi que lorsque $X$ est \emph{réduit}, il en est de même de $\Gamma_u$, qui est isomorphe à $U$ \sref[I]{I.5.3.11} ; alors $\Gamma_\omega$ existe et n'est autre que le sous-préschéma \emph{réduit} de $X\times_S Y$ ayant pour espace sous-jacent l'adhérence dans $X\times_S Y$ de l'espace sous-jacent à $\Gamma_u$ (\sref[I]{I.5.2.1} et \sref[I]{I.5.2.2}).
Notons enfin que si $Y$ est \emph{séparé} sur $S$, $\Gamma_u$ est \emph{fermé} dans $U\times_S Y$ \sref[I]{I.5.4.3}, donc est \emph{induit par} $\Gamma_\omega$ (lorsque ce dernier existe) sur l'ouvert $\Gamma_\omega\cap(U\times_S Y)$ \sref[I]{I.9.5.10} ; par contre, ce préschéma induit est en général différent de $\Gamma_u$ lorsque $Y$ n'est pas séparé.
En particulier, si $v:X\to Y$ est un $S$-morphisme, le graphe $\Gamma_\omega$ de la classe $\omega$ de $v$ peut être distinct du graphe $\Gamma_v$.
Aussi allons-nous dans ce qui suit, lorsqu'il sera question du graphe d'un pseudo-$S$-morphisme, nous borner au cas où \emph{$Y$ est séparé sur $S$}.
\end{env}

\begin{env}[20.4.2]
\label{IV.20.4.2}
Supposons donc $\Gamma_\omega$ défini et $Y$ séparé sur $S$ ; notons $p$ et $q$ les restrictions à $\Gamma_\omega$ des projections canoniques
\[
  \xymatrix{
  & \Gamma_\omega \ar[dl]_p \ar[dr]^q & \\
  X & & Y
  }
\]
Alors, si $U\subset\operatorname{dom}_S(\omega)$, la restriction $p^{-1}(U)\to U$ de $p$ est un isomorphisme \sref[I]{I.5.3.11} ; réciproquement, si $U$ est un ouvert de $X$ ayant cette propriété, et si $u$ est l'isomorphisme réciproque de la restriction $p^{-1}(U)\to U$ de $p$, $q\circ u$ est un $S$-morphisme de $U$ dans $Y$ qui coïncide avec un morphisme de la classe $\omega$ dans $U\cap\operatorname{dom}_S(\omega)$.
On en conclut que $\operatorname{dom}_S(\omega)$ est le \emph{plus grand ouvert} $U$ de $X$ tel que la restriction $p^{-1}(U)\to U$ de $p$ soit un
